\section{Why Clean Sufficiency Is Not Enough}
\label{sec:theory}

The following propositions are information-structure statements.  They explain
why the method must inspect aligned evidence orbits before making item-level
selective RAG decisions.

\paragraph{Proposition 1 (Clean sufficiency does not imply orbit sufficiency).}
\label{prop:clean-not-orbit}
For any clean-only selector that depends only on
$C_i=s_\theta(x_i,a_i,E_i^0)$, there exist two instances with the same clean
score and different orbit risks.  Therefore clean sufficiency alone cannot
identify item-level counterfactual risk.

\paragraph{Argument.}
Construct two items with identical $(x,a,E^0)$ and identical verifier score
$C$.  In the first item, every counterfactual evidence set remains answerable:
$Y^j=1$ for all $j$.  In the second item, the clean set is answerable but one
counterfactual set removes or contradicts the needed support, so $Y^k=0$ for
some $k>0$.  A clean-only selector receives the same observation for both items
and must make the same decision, while their orbit risks differ.  The gap is not
caused by a weak clean verifier; it is caused by missing counterfactual
information.

\paragraph{Proposition 2 (Single-set sufficiency is insufficient).}
\label{prop:single-set-not-orbit}
Any selector that observes one evidence set from an orbit can be made unable to
distinguish a stable orbit from an unstable orbit that agrees on that observed
set.  Orbit-level selective risk is therefore not identifiable from a single
local sufficiency score in general.

\paragraph{Argument.}
Fix the observed set $E^j$ and its score
$s_\theta(x,a,E^j)$.  Create two orbits that share this set and score.  Let the
remaining sets in the first orbit preserve answerability, while at least one
remaining set in the second orbit removes, swaps, or contradicts the support
needed for $a$.  The selector has identical input in both cases but the empirical
orbit risks are different.  This establishes the need for multi-set orbit
statistics such as worst-case sufficiency, dispersion, and clean-to-worst gaps.

\paragraph{Proposition 3 (Orbit alignment is necessary for item-level risk).}
\label{prop:orbit-alignment-necessary}
If counterfactual evidence sets are not aligned to the same query--answer item,
then item-level orbit risk $R_i$ is not identifiable from the observed multiset
of verifier scores and labels.  Alignment is therefore a structural requirement
for CSRM-style risk estimation, not just an implementation detail.

\paragraph{Argument.}
Consider two query--answer items with the same global multiset of
counterfactual evidence scores and labels.  If the assignment of those
counterfactual sets to items is hidden or arbitrarily permuted, the same global
observations can produce two different item-level risk vectors: one assignment
puts the high-risk perturbations on item 1, while another assignment puts them
on item 2.  Both assignments match the global multiset.  Since a selector must
rank or abstain per item, the item-level risk target is underdetermined without
the alignment map from each perturbation back to its clean query--answer frame.

\paragraph{Consequence for CSRM-RAG.}
The propositions justify the method's minimum evidence interface:
$(x_i,a_i,\mathcal{O}_i)$ plus aligned per-set sufficiency scores.  Clean-only
controls, SURE-style single-set controls, and unaligned orbit averages are useful
baselines, but they cannot by themselves establish the counterfactual stability
claim targeted here.

\paragraph{Scope and non-claims.}
These propositions do not assert that every CSRM implementation dominates every
strong baseline, and they do not establish a formal risk-control guarantee.
They only show that the intended claim is impossible to support from clean-only,
single-set, or unaligned evidence information alone.
